%% style: academic
%% desc: 学术笔记风，适合技术分享/论文解读/深度内容 —— 严谨结构化、节标题、要点框
%% 用法：AGENT 读此模板，以 MD 底稿为信息源、按此结构生成 note.tex（不直接替换占位符）。
%% 占位符 %%TITLE%% / %%CONTENT%% 标出内容插入点。
\documentclass[11pt]{article}
\usepackage[UTF8, scheme=plain]{ctex}
\usepackage[margin=2.5cm]{geometry}
\usepackage{amsmath, amssymb}
\usepackage{enumitem}
\usepackage{tcolorbox}
\usepackage{booktabs}
\usepackage{hyperref}
\hypersetup{hidelinks}

\title{%%TITLE%%}
\author{AI 笔记}
\date{\today}

\begin{document}
\maketitle

\begin{abstract}
%%CONTENT%%（摘要：用 3-5 句话概括底稿核心论点）
\end{abstract}

\tableofcontents

% ===== 内容区域：把 MD 底稿按下面结构组织 =====
% 每章用 \section{}，小节用 \subsection{}；要点用 \begin{enumerate}；
% 关键结论/公式用 tcolorbox 强调；术语/定义用 \emph{}。

\section{核心概念}
%% 从底稿提炼定义与关键术语，逐条列出

\section{主要论点}
%% 按逻辑顺序展开底稿论点，每点配解释与例子

\section{关键细节与证据}
%% 底稿中的具体数据、引用、例子，用表格或列表呈现

\section{结论与延伸}
%% 总结 + 与已知知识的联系/延伸思考

\end{document}
